% 对在线学习算法的汇总内容，包括了：
% E-prop, 
% OTTT, NDOT, S-TLLR, TESS
\documentclass[journal,12pt,onecolumn]{IEEEtran}
\IEEEoverridecommandlockouts
\usepackage{cite}
\usepackage{amsmath, amssymb, bm, color, graphicx, enumitem}
\usepackage{amsmath,amssymb,amsfonts}
\usepackage{algorithmic}
\usepackage{algorithm}
%\usepackage{algpseudocode}
\usepackage{graphicx}

\usepackage{booktabs} 

\usepackage{textcomp}
%\usepackage[marginal]{footmisc}
\usepackage{xcolor}
\usepackage{subfigure}
\graphicspath{{./img/}}
\usepackage{tabularx}
\usepackage{makecell}
\usepackage{caption2}
%\usepackage{fontspec}
%\setmainfont{TeX Gyre Termes}
\usepackage{color}

\usepackage[UTF8]{ctex}

\input{math_commands_TESS.tex}
\input{math_commands_STLLR.tex}

\definecolor{myred}{RGB}{255, 67, 20}
\definecolor{myblue}{RGB}{51, 51, 255}
\definecolor{mygreen}{RGB}{0, 128, 0}

\begin{document}
\title{SNN在线与本地学习算法 研究汇总}
\author{Tao}
\date{\today}
\maketitle

\begin{abstract}
SNN在线与本地学习算法 研究汇总
\end{abstract}

\section*{文章汇总}

RTRL: \textbf{A learning algorithm for continually running fully recurrent neural networks} \cite{williams1989learning};

E-prop: \textbf{A solution to the learning dilemma for recurrent networks of spiking neurons} \cite{bellec_solution_2020};

OTTT: \textbf{Online Training Through Time for Spiking Neural Networks} \cite{xiao_online_2022};

NDOT: \textbf{NDOT: Neuronal Dynamics-based Online Training for Spiking Neural Networks} \cite{jiang_ndot_nodate};

S-TLLR: \textbf{S-TLLR: STDP-inspired Temporal Local Learning Rule for Spiking Neural Networks} \cite{apolinario_s-tllr_2024};

TESS: \textbf{TESS: A Scalable Temporally and Spatially Local Learning Rule for Spiking Neural Networks} \cite{apolinario_tess_2025}


\section{Introduction：研究背景与问题引入}

\begin{itemize}
  \item 1. SNN 与 ANN 的基本差异
  \item 2. SNN 训练的核心困难：脉冲的不可求导性质
  \item 3. 时间依赖性与 credit assignment 问题
  \item 4. 传统 BPTT/STBP 的基本思想与使用的局限性
  \item 5. 在线学习与本地学习的研究动机
\end{itemize}

\subsection{SNN 与 ANN 的基本差异}

ANN 使用连续值激活函数：
  \[
    y = f(Wx)
  \]

SNN 使用事件驱动 spike：
  \[
    s_t = H(u_t - \theta)
  \]





在 ANN 中，神经元输出通常是一个连续函数，例如 ReLU 或 sigmoid，输出只依赖当前输入。

但在 SNN 中，神经元的输出是 spike，即一个离散的0-1事件，由膜电位是否超过阈值决定。

更关键的是，SNN 是一个带状态的动态系统，神经元的膜电位会随时间累积并衰减，因此当前时刻的输出不仅依赖当前输入，还依赖历史输入序列。

这一点使得 SNN 更接近生物神经系统，但也使得训练变得更加复杂。


\subsection{SNN 训练的核心困难（不可导性）}


Spike generation function：
  \[
    s_t = H(u_t - \theta)
  \]

Heaviside function：
  \[
    H(x) =
    \begin{cases}
    1 & x > 0 \\
    0 & x \le 0
    \end{cases}
  \]

SNN 训练的核心数学障碍。

SNN 的 spike 是由 Heaviside step function 产生的，它是一个离散函数。

这个函数的导数在绝大多数位置都是零，在阈值点又是不可导的，因此无法直接用于梯度下降。

而深度学习的核心优化方法是 gradient descent梯度下降，因此这构成了 SNN 训练的第一大障碍。

也就是说，我们无法像 ANN 一样直接对 loss function 进行反向传播。



\subsection{时间依赖性与 credit assignment 问题}



\begin{itemize}
  \item SNN neuron state：
  \[
    u_t = f(u_{t-1}, s_{t-1}, x_t)
  \]

  \item 当前输出依赖：
  \begin{itemize}
    \item 当前输入 $x_t$
    \item 历史膜电位 $u_{t-1}$
    \item 历史 spike $s_{t-1}$
  \end{itemize}

  \item 问题：credit assignment over time

  \item 核心困难：
  \begin{itemize}
    \item 哪个时间步导致当前 loss？
    \item 如何将误差分配到过去？
  \end{itemize}
\end{itemize}

SNN 的第二个核心困难：时间依赖性。

在 SNN 中，神经元是一个递归动态系统，它的状态由上一时刻的状态决定。

因此当前的输出不仅由当前输入决定，还受到过去所有时间步的影响。

这导致一个关键问题：当我们计算当前 loss 时，无法直接判断是哪一个时间步的输入或状态导致了这个误差。

这就是 temporal credit assignment problem，也就是时间维度上的误差分配问题。

这个问题是后续所有方法（BPTT、RTRL、E-prop 等）的核心出发点。


\subsection{BPTT / STBP 的基本思想与局限}


\begin{itemize}
  \item 基本思想：
  \begin{itemize}
    \item 将 SNN 沿时间展开
    \item 转换为 deep feedforward graph
  \end{itemize}

  \item STBP：
  \begin{itemize}
    \item surrogate gradient 解决 spike 不可导问题
    \item 结合 BPTT 进行训练
  \end{itemize}

  \item 优点：
  \begin{itemize}
    \item 可以直接使用 gradient descent
    \item 训练效果较好
  \end{itemize}

  \item 局限：
  \begin{itemize}
    \item 需要存储完整时间序列
    \item 非在线更新
    \item 不适合硬件实现
  \end{itemize}
\end{itemize}



\textbf{这一页介绍目前最标准的解决方案：BPTT 和 STBP。}

BPTT 的核心思想是把 SNN 在时间维度展开成一个深度神经网络，然后使用链式法则进行反向传播。

但由于 spike 不可导的问题，需要引入 STBP，也就是 surrogate gradient 方法，用一个平滑函数近似 spike 的导数，从而使梯度可以传播。

这种方法的优点是训练效果好，并且可以直接使用现有的深度学习优化框架。

但它的核心问题是必须存储整个时间序列的状态，因此内存消耗随时间增长，并且无法进行实时在线更新。



\subsection{研究动机：在线学习与本地学习}

\begin{itemize}
  \item BPTT / STBP 的问题：
  \begin{itemize}
    \item memory cost scales with time
    \item not suitable for streaming input
    \item not biologically plausible
  \end{itemize}

  \item 目标1：Online learning
  \begin{itemize}
    \item update weights at each timestep
    \item no need to store full history
  \end{itemize}

  \item 目标2：Local learning
  \begin{itemize}
    \item weight update uses only local information
    \item no global backpropagation required
  \end{itemize}

  \item 本报告主线：
  \[
  \text{BPTT/STBP} \rightarrow \text{Online Learning} \rightarrow \text{Local Learning}
  \]
\end{itemize}


\textbf{这一页总结整个 Introduction 的研究动机。}

虽然 BPTT 和 STBP 可以有效训练 SNN，但它们存在明显的局限性，包括内存消耗随时间增长、无法进行在线更新，以及与生物神经系统机制不一致。

因此研究主要发展出两个方向。

第一个方向是 online learning，也就是在每个时间步直接更新参数，而不依赖完整历史。

第二个方向是 local learning，也就是权重更新只依赖突触局部可获得的信息，而不需要全局误差信号传播。

本报告的后续内容将围绕这一条演化路径展开，从 RTRL E-prop 出发，然后描述 BPTT 和 STBP，逐步介绍在线学习方法以及局部学习规则。




\section{RTRL：实时循环学习算法}

\subsection*{本部分内容概览}
\begin{itemize}
    \item 网络模型与符号约定
    \item 误差定义与目标函数
    \item 核心思想：追踪每个权重对每个输出的影响
    \item 递推公式的详细推导
    \item 权重更新规则与在线版本
    \item 教师强制（Teacher Forcing）变体
    \item 计算复杂度分析
\end{itemize}

\subsubsection*{PPT 内容}
\begin{itemize}
    \item \textbf{网络模型}：$n$ 个单元（神经元），$m$ 个外部输入。
        \begin{itemize}
            \item $\mathbf{y}(t)$：$n$ 维网络输出向量（时刻 $t$）
            \item $\mathbf{x}(t)$：$m$ 维外部输入向量
            \item $\mathbf{z}(t) = [\mathbf{x}(t); \mathbf{y}(t)]$：扩展状态
            \item $U$：网络单元索引集合；$I$：输入索引集合
        \end{itemize}
    \item \textbf{动力学}：
        \begin{equation}
            s_k(t) = \sum_{l \in U \cup I} w_{kl} z_l(t), \quad k \in U
        \end{equation}
        \begin{equation}
            y_k(t+1) = f_k(s_k(t))
        \end{equation}
    \item \textbf{误差}：目标 $d_k(t)$，误差 $e_k(t) = d_k(t) - y_k(t)$ 若 $k\in T(t)$，否则 0。
    \item \textbf{目标}：最小化 $J_{\text{total}} = \sum_{t} J(t)$，$J(t) = \frac{1}{2}\sum_k e_k^2(t)$。
    \item \textbf{核心量}：$p_{ij}^k(t) \triangleq \frac{\partial y_k(t)}{\partial w_{ij}}$。
    \item \textbf{递推}：
        \begin{equation}
            p_{ij}^k(t+1) = f_k'(s_k(t)) \left[ \sum_{l\in U} w_{kl} p_{ij}^l(t) + \delta_{ik} z_j(t) \right]
        \end{equation}
        初始条件 $p_{ij}^k(t_0)=0$。
    \item \textbf{权重更新}：$\Delta w_{ij}(t) = \alpha \sum_{k\in U} e_k(t) p_{ij}^k(t)$（在线版本）。
    \item \textbf{教师强制}：用目标值替代实际输出，递推中排除被强制单元。
    \item \textbf{复杂度}：$O(n^3 + mn^2)$ 存储，$O(n^4)$ 计算每步。
\end{itemize}

\subsubsection*{讲稿（详细推导）}

\paragraph{1. 网络模型与符号}
我们考虑一个完全循环的神经网络，共有 $n$ 个单元（神经元）和 $m$ 个外部输入线。在离散时间 $t$，网络输出向量记为 $\mathbf{y}(t) \in \mathbb{R}^n$，外部输入向量记为 $\mathbf{x}(t) \in \mathbb{R}^m$。将两者拼接得到扩展向量 $\mathbf{z}(t) = [\mathbf{x}(t); \mathbf{y}(t)] \in \mathbb{R}^{m+n}$。我们定义

$U$ 为单元索引集合（$|U|=n$）;

$I$ 为输入索引集合（$|I|=m$）;

为了统一处理，给每个输入也分配一个索引，

使得 $z_k(t)$ 当 $k\in I$ 时取 $x_k(t)$，当 $k\in U$ 时取 $y_k(t)$。

网络中的每个单元 $k$ 在时刻 $t$ 接收到净输入 $s_k(t)$，其计算方式为所有扩展输入 $z_l(t)$ 乘上对应权重 $w_{kl}$ 之和：
\[
s_k(t) = \sum_{l\in U \cup I} w_{kl} z_l(t).
\]
然后单元的输出由非线性激活函数 $f_k$ 作用，在下一时刻产生：
\[
y_k(t+1) = f_k(s_k(t)).
\]
这里 $f_k$ 通常取 sigmoid 或 tanh 函数。注意，外部输入 $x(t)$ 并不会立刻影响输出，而是延迟一个时间步。

整个网络的状态由所有单元的膜电位（或激活值）和输入共同决定，权重矩阵 $\mathbf{W}$ 是一个 $n\times(m+n)$ 的矩阵，包含了所有单元间和输入到单元的连接。为了引入偏置，我们可以在输入中包含一个始终为 1 的常数输入。

\paragraph{2. 误差与目标函数}
在监督学习场景中，某些单元在特定时刻被指定为目标值。令:

$T(t)$ 为在时刻 $t$ 有目标值的单元集合;

目标值为 $d_k(t)$;

定义误差 $e_k(t)$：

\[
e_k(t) =
\begin{cases}
d_k(t) - y_k(t), & k\in T(t) \\
0, & \text{否则}.
\end{cases}
\]
瞬时误差 $J(t) = \frac{1}{2}\sum_{k\in U} e_k^2(t)$，总误差 $J_{\text{total}} = \sum_{t=t_0+1}^{t_1} J(t)$。我们的目标是通过梯度下降调整权重 $\mathbf{W}$，使 $J_{\text{total}}$ 最小。

\paragraph{3. 核心思想：跟踪敏感度 \(p_{ij}^k\)}

资格迹：

跟踪敏感度 \(p_{ij}^k\)


RTRL 的关键创新是定义一个三阶量 \(p_{ij}^k(t)\)，它表示时刻 \(t\) 单元 \(k\) 的输出 \(y_k(t)\) 对权重 \(w_{ij}\) 的偏导数：
\[
p_{ij}^k(t) \triangleq \frac{\partial y_k(t)}{\partial w_{ij}}.
\]
这个量回答了这样一个问题：如果我们把权重 \(w_{ij}\) 微调一点点，单元 \(k\) 在时刻 \(t\) 的输出会变化多少？

由于网络是循环的，\(y_k(t)\) 不仅直接依赖于某些权重（如果 \(k=i\)，即权重 \(w_{ij}\) 直接连到单元 \(i\) 时），还通过循环连接间接依赖于所有权重——因为过去的输出会影响当前的膜电位，而膜电位又决定当前的输出。因此，\(p_{ij}^k(t)\) 必须递归计算。

初始时刻 \(t_0\)，我们假设网络的初始状态（如初始膜电位）与所有权重无关，因此：
\[
p_{ij}^k(t_0) = 0, \quad \forall i,j,k.
\]

\paragraph{4. 递推公式的详细推导（分步拆解）}

--


我们的目标是推导 \(p_{ij}^k(t+1)\) 的递推式。由于 \(y_k(t+1) = f_k(s_k(t))\)，应用链式法则，第一步总是：
\[
p_{ij}^k(t+1) = f_k'(s_k(t)) \cdot \frac{\partial s_k(t)}{\partial w_{ij}}. \tag{4.1}
\]
因此，核心问题转化为：计算净输入 \(s_k(t)\) 对权重 \(w_{ij}\) 的偏导数 \(\frac{\partial s_k(t)}{\partial w_{ij}}\)。

\paragraph{4.1 展开 \(s_k(t)\) 的表达式}

净输入 \(s_k(t)\) 是所有来源信号乘以对应权重的总和：
\[
s_k(t) = \sum_{l \in U \cup I} w_{kl} \, z_l(t). \tag{4.2}
\]
其中 \(U\) 是网络单元的索引集合，\(I\) 是外部输入的索引集合。扩展向量 \(z_l(t)\) 定义为：
\[
z_l(t) = 
\begin{cases}
y_l(t), & l \in U, \\
x_l(t), & l \in I.
\end{cases}
\]

为了更清楚地看到权重对输出的依赖性，我们可以把 (4.2) 拆成两部分：
\[
s_k(t) = \sum_{l \in U} w_{kl} \, y_l(t) + \sum_{l \in I} w_{kl} \, x_l(t). \tag{4.3}
\]
这里的第二项（输入部分）不会引入对权重的隐式依赖，因为 \(x_l(t)\) 与权重无关。

\paragraph{4.2 对 (4.3) 直接求导——显式依赖与隐式依赖}

现在对 (4.3) 关于 \(w_{ij}\) 求导。注意：\(y_l(t)\) 本身也是所有权重的函数（因为过去的权重影响了过去的输出），所以必须区分两类依赖：

\begin{enumerate}
    \item \textbf{显式依赖（直接）}：\(s_k(t)\) 的表达式里，哪些地方明确写了 \(w_{ij}\)？只有那些系数恰好是 \(w_{ij}\) 的项。
    \item \textbf{隐式依赖（间接）}：\(y_l(t)\) 包含 \(w_{ij}\)，所以 \(\sum w_{kl} y_l(t)\) 这一整项也会随 \(w_{ij}\) 变化。
\end{enumerate}

因此，我们对 (4.3) 求导，得到三项：
\[
\frac{\partial s_k(t)}{\partial w_{ij}} = 
\underbrace{\sum_{l \in U} w_{kl} \frac{\partial y_l(t)}{\partial w_{ij}}}_{\text{隐式依赖：通过循环输出传播}}
+
\underbrace{\sum_{l \in U} \frac{\partial w_{kl}}{\partial w_{ij}} y_l(t)}_{\text{显式依赖：来自单元连接的系数}}
+
\underbrace{\sum_{l \in I} \frac{\partial w_{kl}}{\partial w_{ij}} x_l(t)}_{\text{显式依赖：来自输入连接的系数}}. \tag{4.4}
\]

\paragraph{4.3 逐一计算三项的具体值}

\textbf{第一项（隐式依赖）}：
\[
\frac{\partial y_l(t)}{\partial w_{ij}} = p_{ij}^l(t),
\]
这正是我们在前一个时间步已经定义好的量。所以第一项为：
\[
\sum_{l \in U} w_{kl} \, p_{ij}^l(t). \tag{4.5}
\]

\textbf{第二项（显式依赖，来自单元连接）}：
我们问：在 \(\sum_{l \in U} w_{kl} y_l(t)\) 中，哪一项的系数恰好是我们要微调的 \(w_{ij}\)？

- 项的通式是 \(w_{kl} \cdot y_l(t)\)。

- 这个系数 \(w_{kl}\) 等于 \(w_{ij}\) 的条件是：\textbf{\(k = i\) 且 \(l = j\)}。

- 因为 \(w_{ij}\) 的下标中，\(i\) 是目标神经元，\(j\) 是源头（源头可能是单元）。

- 所以只有当 \(k=i\)（我们正在计算净输入的神经元恰好是权重所在的目标神经元）且 \(l=j\)（这个项恰好来自源头 \(j\)）时，导数为 \(y_j(t)\)。

- 其他情况下，导数为 0。

因此：
\[
\sum_{l \in U} \frac{\partial w_{kl}}{\partial w_{ij}} y_l(t) = 
\begin{cases}
y_j(t), & \text{如果 } k=i \text{ 且 } j \in U,\\
0, & \text{否则}.
\end{cases} \tag{4.6}
\]

\textbf{第三项（显式依赖，来自输入连接）}：
同理，在 \(\sum_{l \in I} w_{kl} x_l(t)\) 中，系数 \(w_{kl}\) 等于 \(w_{ij}\) 的条件是 \textbf{\(k=i\) 且 \(l=j\)}，并且此时 \(j\) 必须是输入（\(j \in I\)）。导数为 \(x_j(t)\)。

因此：
\[
\sum_{l \in I} \frac{\partial w_{kl}}{\partial w_{ij}} x_l(t) = 
\begin{cases}
x_j(t), & \text{如果 } k=i \text{ 且 } j \in I,\\
0, & \text{否则}.
\end{cases} \tag{4.7}
\]

\paragraph{4.4 引入克罗内克函数 \(\delta_{ik}\)}

为了简洁地写出 (4.6) 和 (4.7) 的共同规律，我们引入克罗内克函数 \(\delta_{ik}\)，其定义为：
\[
\delta_{ik} = 
\begin{cases}
1, & \text{如果 } i = k,\\
0, & \text{如果 } i \neq k.
\end{cases}
\]

这个函数的作用就像一个“开关”或“门锁”——它只在 \(i=k\) 时打开，其他时候保持关闭。

有了这个记号，我们可以把第二项和第三项合并为：
\[
\delta_{ik} \cdot z_j(t). \tag{4.8}
\]

为什么可以这样合并？因为：
- 如果 \(j \in U\)，则 \(z_j(t) = y_j(t)\)，且第二项非零。
- 如果 \(j \in I\)，则 \(z_j(t) = x_j(t)\)，且第三项非零。
- 无论哪种情况，其值恰好是 \(\delta_{ik} \cdot z_j(t)\)。

因此，(4.4) 的三项可以合并为：
\[
\frac{\partial s_k(t)}{\partial w_{ij}} = \sum_{l \in U} w_{kl} \, p_{ij}^l(t) + \delta_{ik} \, z_j(t). \tag{4.9}
\]

\paragraph{4.5 如何理解 \(\delta_{ik} z_j(t)\) 这一项？}

这一项的含义非常直观：
\begin{itemize}
    \item \textbf{只有当你正在计算的那个神经元 \(k\) 恰好是权重 \(w_{ij}\) 所在的神经元 \(i\) 时}（即 \(k=i\)），直接项 \(z_j(t)\) 才会出现。
    \item 如果 \(k \neq i\)，则 \(s_k(t)\) 的表达式中根本不包含 \(w_{ij}\) 这个系数，所以直接项为 0。
    \item 而 \(\sum_{l \in U} w_{kl} p_{ij}^l(t)\) 这一项对所有 \(k\) 都存在，因为循环连接使得所有单元的历史输出都间接依赖于 \(w_{ij}\)。
\end{itemize}

所以，完整的影响路径是：
\[
\boxed{\text{历史影响（通过循环传播，所有人都能感受到）} + \text{直接影响（仅权重所在的神经元能直接感受到）}.}
\]

\paragraph{4.6 代回得到最终递推式}

将 (4.9) 代入 (4.1)：
\[
\boxed{
p_{ij}^k(t+1) = f_k'(s_k(t)) \left[ \sum_{l \in U} w_{kl} \, p_{ij}^l(t) + \delta_{ik} \, z_j(t) \right].
} \tag{4.10}
\]
这就是 RTRL 的核心递推方程，它告诉我们如何从前一个时刻的 \(p_{ij}^l(t)\) 计算出下一个时刻的 \(p_{ij}^k(t+1)\)。这里 \(j\) 可以属于 \(U\)（单元）或 \(I\)（输入），而 \(i \in U\)，\(k \in U\)。

\paragraph{5. 权重更新规则}

有了 \(p_{ij}^k(t)\) 之后，我们可以计算损失 \(J(t) = \frac{1}{2}\sum_k e_k^2(t)\) 对权重 \(w_{ij}\) 的梯度：
\[
-\frac{\partial J(t)}{\partial w_{ij}} = \sum_{k \in U} e_k(t) \frac{\partial y_k(t)}{\partial w_{ij}} = \sum_{k \in U} e_k(t) \, p_{ij}^k(t). \tag{5.1}
\]
其中误差 \(e_k(t)\) 在之前已定义。

\paragraph{5.1 离线累积版本}
如果固定权重跑完整个轨迹再更新，则累积所有时刻的梯度：

离线累积：

\[
\Delta w_{ij} = -\alpha \sum_{t} \frac{\partial J(t)}{\partial w_{ij}} = \alpha \sum_{t} \sum_{k} e_k(t) \, p_{ij}^k(t). \tag{5.2}
\]

\paragraph{5.2 在线实时版本}
如果边跑边更新，则每个时刻独立更新：

在线实时：
\[
\Delta w_{ij}(t) = \alpha \sum_{k} e_k(t) \, p_{ij}^k(t), \quad w_{ij} := w_{ij} + \Delta w_{ij}(t). \tag{5.3}
\]
这种在线版本不需要定义训练区间，网络可以持续运行并学习。

\paragraph{6. 教师强制（Teacher Forcing）变体}
---


在某些任务中，我们可以把单元 \(k\) 在时刻 \(t\) 的输出替换为目标值 \(d_k(t)\)，而不是使用它实际产生的输出 \(y_k(t)\)。这种技巧称为教师强制。

在教师强制下，\(z_k(t)\) 的定义变为：
\[
z_k(t) = 
\begin{cases}
d_k(t), & \text{如果 } k \in T(t),\\
y_k(t), & \text{如果 } k \in U \setminus T(t),
\end{cases}
\]
其中 \(T(t)\) 是在时刻 \(t\) 有目标值的单元集合。

因为教师值 \(d_k(t)\) 不依赖于任何权重，所以 \(\frac{\partial d_k(t)}{\partial w_{ij}} = 0\)。这意味着在递推式中，被教师强制的单元对应的 \(p_{ij}^l(t)\) 项必须被移除。因此递推式变为：
\[
p_{ij}^k(t+1) = f_k'(s_k(t)) \left[ \sum_{l \in U \setminus T(t)} w_{kl} \, p_{ij}^l(t) + \delta_{ik} \, z_j(t) \right]. \tag{6.1}
\]
这等价于在计算 \(\frac{\partial s_k(t)}{\partial w_{ij}}\) 时，把被教师强制单元的 \(\frac{\partial y_l(t)}{\partial w_{ij}}\) 视为 0。

\paragraph{7. 计算复杂度分析}

对于全连接网络：
\begin{itemize}
    \item 有 \(n\) 个单元，\(m\) 个外部输入。
    \item 权重总数 = \(n \times (m+n)\)。
    \item 对于每个权重，需要存储它对所有 \(n\) 个单元输出的导数 \(p_{ij}^k\)。
\end{itemize}

因此需要存储的变量总数为：
\[
\#p = n \times n \times (m+n) = n^3 + m n^2. \tag{7.1}
\]

每时间步更新每个 \(p_{ij}^k(t+1)\) 需要计算：
\[
\sum_{l \in U} w_{kl} \, p_{ij}^l(t),
\]
这是一个对 \(n\) 项的求和，复杂度为 \(O(n)\)。因此每时间步的总计算复杂度为：
\[
O(n) \times O(n^3) = O(n^4). \tag{7.2}
\]
这使得 RTRL 在大规模网络中极不实用，但它是第一个真正意义上的在线梯度算法，为后续工作（如 E-prop）奠定了理论基础。

\section{E-prop：资格传播}

\subsection*{本部分内容概览}
\begin{itemize}
    \item 脉冲神经网络的一般模型与符号
    \item 核心数学定理：梯度分解为学习信号与资格迹的乘积
    \item 资格迹的定义与递推
    \item 定理的严格证明（从 BPTT 到 E-prop 的重新因子化）
    \item LIF 与 ALIF 神经元的具体资格迹推导
    \item 学习信号的不同变体（对称、随机、自适应）
    \item 奖励基 E-prop（强化学习）
\end{itemize}

\subsubsection*{PPT 内容}
\begin{itemize}
    \item \textbf{通用模型}：可观状态 $z_j^t$（脉冲 0/1），隐藏状态 $\mathbf{h}_j^t$（膜电位、适应变量等）。
        \begin{equation}
            \mathbf{h}_j^t = M(\mathbf{h}_j^{t-1}, z^{t-1}, \mathbf{x}^t, \mathbf{W}_j)
        \end{equation}
        \begin{equation}
            z_j^t = f(\mathbf{h}_j^t)
        \end{equation}
    \item \textbf{核心目标}：计算损失 $E$ 对突触权重 $W_{ji}$ 的梯度 $\frac{dE}{dW_{ji}}$。
    \item \textbf{关键定理}（式 (1)(3)）：
        \begin{equation}
            \frac{dE}{dW_{ji}} = \sum_t L_j^t \, e_{ji}^t
        \end{equation}
        其中 $L_j^t = \frac{dE}{dz_j^t}$（学习信号），$e_{ji}^t = \left[\frac{dz_j^t}{dW_{ji}}\right]_{\text{local}}$（资格迹）。
    \item \textbf{资格迹定义}：
        \begin{equation}
            e_{ji}^t = \frac{\partial z_j^t}{\partial \mathbf{h}_j^t} \cdot \boldsymbol{\epsilon}_{ji}^t
        \end{equation}
        \begin{equation}
            \boldsymbol{\epsilon}_{ji}^t = \sum_{t'\le t} \left( \frac{\partial \mathbf{h}_j^t}{\partial \mathbf{h}_j^{t-1}} \cdots \frac{\partial \mathbf{h}_j^{t'+1}}{\partial \mathbf{h}_j^{t'}} \right) \cdot \frac{\partial \mathbf{h}_j^{t'}}{\partial W_{ji}}
        \end{equation}
        递推：$\boldsymbol{\epsilon}_{ji}^t = \frac{\partial \mathbf{h}_j^t}{\partial \mathbf{h}_j^{t-1}} \cdot \boldsymbol{\epsilon}_{ji}^{t-1} + \frac{\partial \mathbf{h}_j^t}{\partial W_{ji}}$
    \item \textbf{定理证明要点}：从 BPTT 经典因子化出发，递归展开 $\frac{dE}{d\mathbf{h}_j^{t'}}$，交换求和顺序，得到 $L_j^t$ 与 $e_{ji}^t$ 的乘积形式。
    \item \textbf{LIF 神经元的资格迹}：
        \begin{equation}
            e_{ji}^t = \psi_j^t \, \bar{z}_i^{t-1}, \quad \bar{z}_i^{t-1} = \sum_{\tau} \alpha^{t-1-\tau} z_i^\tau
        \end{equation}
    \item \textbf{ALIF 神经元（适应）}：存在额外慢变量 $a_j^t$，资格迹包含慢分量。
        \begin{equation}
            e_{ji}^t = \psi_j^t \left( \bar{z}_i^{t-1} - \beta \epsilon_{ji,a}^t \right)
        \end{equation}
        其中 $\epsilon_{ji,a}^{t+1} = \psi_j^t \bar{z}_i^{t-1} + (\rho - \psi_j^t\beta) \epsilon_{ji,a}^t$
    \item \textbf{学习信号变体}：
        \begin{itemize}
            \item 对称：$L_j^t = \sum_k W_{kj}^{out}(y_k^t - y_k^{*,t})$
            \item 随机：$L_j^t = \sum_k B_{jk}(y_k^t - y_k^{*,t})$（固定随机）
            \item 自适应：随机加可塑性
        \end{itemize}
    \item \textbf{奖励基 E-prop}：用于强化学习，$\Delta W_{ji} = -\eta \sum_t \delta^t \mathcal{F}_\gamma( L_j^t \bar{e}_{ji}^t )$，$\delta^t$ 为奖励预测误差。
\end{itemize}

\subsubsection*{讲稿（详细推导）}

\paragraph{1. 脉冲网络的通用模型}
E-prop 面向脉冲循环神经网络（RSNN）。我们考虑一般形式：每个神经元 $j$ 具有隐藏状态向量 $\mathbf{h}_j^t$（可能包含膜电位、适应变量等）和可观输出 $z_j^t \in \{0,1\}$（脉冲）。网络动力学由两个函数定义：
\[
\mathbf{h}_j^t = M(\mathbf{h}_j^{t-1}, z^{t-1}, \mathbf{x}^t, \mathbf{W}_j), \quad z_j^t = f(\mathbf{h}_j^t).
\]
其中 $M$ 是状态更新函数（例如 LIF 膜电位更新），$f$ 是脉冲生成函数（通常为 Heaviside 阶跃函数）。权重 $W_{ji}$ 是从神经元 $i$ 到 $j$ 的突触权重。

\paragraph{2. 梯度分解定理}
E-prop 的核心数学结果是，损失函数 $E$（取决于所有脉冲）对权重 $W_{ji}$ 的梯度可以表示为：
\[
\frac{dE}{dW_{ji}} = \sum_t L_j^t \, e_{ji}^t,
\]
其中 $L_j^t = \frac{dE}{dz_j^t}$ 称为“学习信号”，$e_{ji}^t$ 称为“资格迹”。资格迹定义为局部可计算的量：
\[
e_{ji}^t \triangleq \left[ \frac{dz_j^t}{dW_{ji}} \right]_{\text{local}}.
\]
这个局部导数不是近似，它收集了所有可以通过前向计算（仅依赖神经元 $j$ 和 $i$ 的历史）获得的梯度信息。具体展开为：
\[
e_{ji}^t = \frac{\partial z_j^t}{\partial \mathbf{h}_j^t} \cdot \boldsymbol{\epsilon}_{ji}^t,
\]
其中 $\boldsymbol{\epsilon}_{ji}^t$ 是“资格向量”，定义为：
\[
\boldsymbol{\epsilon}_{ji}^t = \sum_{t'\le t} \left( \frac{\partial \mathbf{h}_j^t}{\partial \mathbf{h}_j^{t-1}} \cdots \frac{\partial \mathbf{h}_j^{t'+1}}{\partial \mathbf{h}_j^{t'}} \right) \cdot \frac{\partial \mathbf{h}_j^{t'}}{\partial W_{ji}}.
\]
资格向量满足简单的递归方程：
\[
\boldsymbol{\epsilon}_{ji}^t = \frac{\partial \mathbf{h}_j^t}{\partial \mathbf{h}_j^{t-1}} \cdot \boldsymbol{\epsilon}_{ji}^{t-1} + \frac{\partial \mathbf{h}_j^t}{\partial W_{ji}},
\]
初始 $\boldsymbol{\epsilon}_{ji}^{0}=0$。这个递归完全在前向过程中计算，无需存储历史。

\paragraph{3. 定理的严格证明}
我们从 BPTT 的经典因子化开始。将 RNN 按时间展开成前馈网络，每个时间步 $t'$ 对应一层，权重共享。对于前馈网络，损失对权重的梯度为：
\[
\frac{dE}{dW_{ji}} = \sum_{t'} \frac{dE}{d\mathbf{h}_j^{t'}} \cdot \frac{\partial \mathbf{h}_j^{t'}}{\partial W_{ji}}.
\]
这是因为每个时间步的权重副本相同，求和。注意 $\frac{dE}{d\mathbf{h}_j^{t'}}$ 是总导数，考虑了后续时刻的影响。

现在递归展开 $\frac{dE}{d\mathbf{h}_j^{t'}}$。在计算图节点 $\mathbf{h}_j^{t'}$，它连接到 $z_j^{t'}$ 和 $\mathbf{h}_j^{t'+1}$，因此：
\[
\frac{dE}{d\mathbf{h}_j^{t'}} = \frac{dE}{dz_j^{t'}} \frac{\partial z_j^{t'}}{\partial \mathbf{h}_j^{t'}} + \frac{dE}{d\mathbf{h}_j^{t'+1}} \frac{\partial \mathbf{h}_j^{t'+1}}{\partial \mathbf{h}_j^{t'}}.
\]
定义 $L_j^{t'} = \frac{dE}{dz_j^{t'}}$，并继续展开 $\frac{dE}{d\mathbf{h}_j^{t'+1}}$，直到最后时间 $T$（边界为零），得到：
\[
\frac{dE}{d\mathbf{h}_j^{t'}} = \sum_{t\ge t'} L_j^t \frac{\partial z_j^t}{\partial \mathbf{h}_j^t} \left( \frac{\partial \mathbf{h}_j^t}{\partial \mathbf{h}_j^{t-1}} \cdots \frac{\partial \mathbf{h}_j^{t'+1}}{\partial \mathbf{h}_j^{t'}} \right).
\]
将此代入经典因子化：
\[
\frac{dE}{dW_{ji}} = \sum_{t'} \sum_{t\ge t'} L_j^t \frac{\partial z_j^t}{\partial \mathbf{h}_j^t} \left( \prod_{\tau=t'+1}^{t} \frac{\partial \mathbf{h}_j^{\tau}}{\partial \mathbf{h}_j^{\tau-1}} \right) \frac{\partial \mathbf{h}_j^{t'}}{\partial W_{ji}}.
\]
交换求和顺序（先 $t$ 后 $t'$）：
\[
\frac{dE}{dW_{ji}} = \sum_t L_j^t \frac{\partial z_j^t}{\partial \mathbf{h}_j^t} \left[ \sum_{t'\le t} \left( \prod_{\tau=t'+1}^{t} \frac{\partial \mathbf{h}_j^{\tau}}{\partial \mathbf{h}_j^{\tau-1}} \right) \frac{\partial \mathbf{h}_j^{t'}}{\partial W_{ji}} \right].
\]
方括号中的项正是 $\boldsymbol{\epsilon}_{ji}^t$ 的定义，因此得到 $\frac{dE}{dW_{ji}} = \sum_t L_j^t e_{ji}^t$。这完成了证明。

\paragraph{4. LIF 神经元资格迹}
对于 LIF 神经元，隐藏状态仅为膜电位 $v_j^t$，隐藏状态更新为：
\[
v_j^{t+1} = \alpha v_j^t + \sum_i W_{ji} z_i^t - z_j^t v_{\text{th}}.
\]
忽略重置项对导数的贡献（或按公式处理），我们有 $\frac{\partial v_j^{t+1}}{\partial v_j^t} = \alpha$，$\frac{\partial v_j^t}{\partial W_{ji}} = z_i^{t-1}$（近似）。资格向量变为：
\[
\epsilon_{ji}^t = \sum_{t'\le t} \alpha^{t-t'} z_i^{t'-1} = \bar{z}_i^{t-1},
\]
即前突触脉冲的低通滤波。脉冲导数 $\frac{\partial z_j^t}{\partial v_j^t}$ 用伪梯度 $\psi_j^t$ 代替，故 $e_{ji}^t = \psi_j^t \bar{z}_i^{t-1}$。

\paragraph{5. ALIF 神经元资格迹}
ALIF 神经元有适应变量 $a_j^t$，阈值 $A_j^t = v_{\text{th}} + \beta a_j^t$，适应动态 $a_j^{t+1} = \rho a_j^t + z_j^t$。隐藏状态为二维。资格向量有两个分量，其中适应分量满足：
\[
\epsilon_{ji,a}^{t+1} = \psi_j^t \bar{z}_i^{t-1} + (\rho - \psi_j^t \beta) \epsilon_{ji,a}^t.
\]
因为 $\frac{\partial z_j^t}{\partial \mathbf{h}_j^t} = [\psi_j^t, -\beta \psi_j^t]$，最终资格迹为：
\[
e_{ji}^t = \psi_j^t \left( \bar{z}_i^{t-1} - \beta \epsilon_{ji,a}^t \right).
\]
这引入了慢分量，使梯度信息能跨越长时延传播。

\paragraph{6. 学习信号变体}
理想学习信号 $L_j^t = \frac{dE}{dz_j^t}$ 需要反向传播。E-prop 使用近似：
\begin{itemize}
    \item 对称：$L_j^t = \sum_k W_{kj}^{out} (y_k^t - y_k^{*,t})$，其中 $W_{kj}^{out}$ 是输出权重，$y_k^t$ 是输出神经元的滤波活性。这要求权重对称。
    \item 随机：$L_j^t = \sum_k B_{jk} (y_k^t - y_k^{*,t})$，$B_{jk}$ 固定随机，类似反馈对齐。
    \item 自适应：在随机基础上，$B_{jk}$ 通过局部可塑性逐渐演化，模仿对称情况。
\end{itemize}

\paragraph{7. 奖励基 E-prop}
对于强化学习，网络输出策略和价值估计。使用奖励预测误差 $\delta^t = r^t + \gamma V^{t+1} - V^t$。权重更新为：
\[
\Delta W_{ji} = -\eta \sum_t \delta^t \mathcal{F}_\gamma( L_j^t \bar{e}_{ji}^t ),
\]
其中 $\mathcal{F}_\gamma$ 是折扣因子低通滤波，$\bar{e}_{ji}^t$ 是资格迹的平滑版本。这使 E-prop 适用于深度强化学习任务。

\section{RTRL 与 E-prop 的联系与对比}

\subsection*{本部分内容概览}
\begin{itemize}
    \item 数学结构的对比
    \item 复杂度的对比
    \item 生物合理性的对比
    \item 对脉冲神经元的适应性
    \item 总结与演化关系
\end{itemize}

\subsubsection*{PPT 内容}
\begin{itemize}
    \item \textbf{数学结构}：
        \begin{itemize}
            \item RTRL：$\Delta w_{ij}(t) = \alpha \sum_k e_k(t) p_{ij}^k(t)$，$p_{ij}^k$ 为三阶张量。
            \item E-prop：$\Delta W_{ji} = -\eta \sum_t L_j^t e_{ji}^t$，分解为学习信号和资格迹。
            \item E-prop 可视为 RTRL 的因子化：$p_{ij}^k(t)$ 被分解为 $L_j^t$ 和 $e_{ji}^t$ 的乘积（求和）。
        \end{itemize}
    \item \textbf{复杂度}：
        \begin{itemize}
            \item RTRL：存储 $O(n^3)$，计算 $O(n^4)$（全连接）。
            \item E-prop：存储 $O(S)$（突触数），计算 $O(S)$。
        \end{itemize}
    \item \textbf{生物合理性}：
        \begin{itemize}
            \item RTRL：非局部（需要全局误差和权重矩阵），不生物。
            \item E-prop：局部资格迹 + 全局学习信号（类比多巴胺），更生物合理。
        \end{itemize}
    \item \textbf{适应性}：
        \begin{itemize}
            \item RTRL：主要针对连续值单元，不直接适用于脉冲神经元。
            \item E-prop：专为脉冲神经元设计，包含适应神经元的慢变量。
        \end{itemize}
    \item \textbf{总结}：E-prop 是 RTRL 在脉冲网络中的高效、局部、生物合理的演化版本。
\end{itemize}

\subsubsection*{讲稿}

\paragraph{1. 数学结构的对比}
RTRL 的核心是 $p_{ij}^k(t) = \frac{\partial y_k(t)}{\partial w_{ij}}$，它是一个三阶张量，需要追踪每个权重对每个神经元输出的影响。权重更新需要累加所有输出单元的误差 $e_k(t)$ 与相应 $p_{ij}^k$ 的乘积。而 E-prop 通过分解将这种影响分成两部分：学习信号 $L_j^t$（仅依赖于神经元 $j$ 的全局影响）和资格迹 $e_{ji}^t$（仅依赖突触前后的局部活动）。这样，原本的 $p_{ij}^k$ 被替换为 $L_j^t \cdot e_{ji}^t$ 在时间上的求和，大大简化了存储和计算。

\paragraph{2. 复杂度的对比}
RTRL 需要存储 $n^3$ 量级的变量（全连接），计算每步 $O(n^4)$，无法扩展。E-prop 只需存储每个突触的资格迹（$O(S)$，$S$ 为突触数），计算也仅与突触数线性相关，这是最优的，因为模拟网络本身就需要 $O(S)$ 时间。

\paragraph{3. 生物合理性}
RTRL 的更新需要每个突触访问所有单元的误差和权重矩阵，这在大脑中无法实现。E-prop 的资格迹仅依赖于前突触和后突触的局部活动（如脉冲频率和膜电位），而学习信号可以是全局的（如多巴胺释放），这在生物学上有对应。因此 E-prop 更符合神经可塑性的约束。

\paragraph{4. 对脉冲神经元的适应性}
RTRL 最初设计用于连续激活函数，而脉冲神经元的离散性需要额外的处理（如伪导数）。E-prop 从头针对脉冲设计，并且自然地包含了 ALIF 神经元的慢适应变量，使其能够解决长时间信用分配问题。BPTT 也能处理脉冲网络，但 RTRL 直接用于脉冲网络会面临更大的计算和导数问题。

\paragraph{5. 总结}
E-prop 可以视为 RTRL 在脉冲网络中的高效、局部化、生物合理的演化版本。它继承了 RTRL 的在线学习能力，但通过重新因子化梯度大幅降低了复杂度，并适配了脉冲神经元的特性。两者共同为循环网络的在线学习提供了理论基础，其中 E-prop 更加实用和生物可信。







\section{脉冲神经网络反向传播算法：BPTT}


\begin{equation}
\small
    \frac{\partial L}{\partial \mathbf{W}^l}=\sum_{t=1}^T \frac{\partial L}{\partial \mathbf{s}^{l+1}[t]}{\color{red}\frac{\partial \mathbf{s}^{l+1}[t]}{\partial \mathbf{u}^{l+1}[t]}}\left(\frac{\partial \mathbf{u}^{l+1}[t]}{\partial \mathbf{W}^l}+\sum_{\tau < t}\prod_{i=t-1}^{\tau}\left({\color[rgb]{0,0.69,0.31}\frac{\partial \mathbf{u}^{l+1}[i+1]}{\partial \mathbf{u}^{l+1}[i]}}+{\color{cyan}\frac{\partial \mathbf{u}^{l+1}[i+1]}{\partial \mathbf{s}^{l+1}[i]}}{\color{red}\frac{\partial \mathbf{s}^{l+1}[i]}{\partial \mathbf{u}^{l+1}[i]}}\right)\frac{\partial \mathbf{u}^{l+1}[\tau]}{\partial \mathbf{W}^l}\right),
    \label{bptt gradient}
\end{equation} \cite{xiao_online_2022};

\begin{equation}
{\color[rgb]{0,0.69,0.31}\frac{\partial \mathbf{u}^{l+1}[i+1]}{\partial \mathbf{u}^{l+1}[i]}}+{\color{cyan}\frac{\partial \mathbf{u}^{l+1}[i+1]}{\partial \mathbf{s}^{l+1}[i]}}{\color{red}\frac{\partial \mathbf{s}^{l+1}[i]}{\partial \mathbf{u}^{l+1}[i]}}
\end{equation}

\[
\boxed{
\frac{d \mathbf{u}^{l+1}[i+1]}{d \mathbf{u}^{l+1}[i]}
=
{\color[rgb]{0,0.69,0.31}\frac{\partial \mathbf{u}^{l+1}[i+1]}{\partial \mathbf{u}^{l+1}[i]}}
+
{\color{cyan}\frac{\partial \mathbf{u}^{l+1}[i+1]}{\partial \mathbf{s}^{l+1}[i]}}
\cdot
{\color{red}\frac{\partial \mathbf{s}^{l+1}[i]}{\partial \mathbf{u}^{l+1}[i]}}
}
\]

代入 LIF 动力学后展开为：
\[
\boxed{
\frac{d \mathbf{u}^{l+1}[i+1]}{d \mathbf{u}^{l+1}[i]}
= \lambda \mathbf{I} - \lambda V_{th} \mathbf{I} \cdot \sigma'(\mathbf{u}^{l+1}[i] - V_{th})
}
\]
% ------------------------------------------------------------
% 内容概览（导航）
% ------------------------------------------------------------
\subsection*{本部分核心结构}
\begin{itemize}[leftmargin=*, itemsep=0.2ex]
    \item[1.] 前向时空动态与LIF离散方程
    \item[2.] BPTT完整梯度公式（含颜色标记）
    \item[3.] 公式拆解：误差传入、当前直连项
    \item[4.] 公式拆解：时间传播核（泄漏 vs 重置路径）
    \item[5.] 代理梯度（SG）、STBP严格定义及BPTT固有瓶颈
\end{itemize}
\vspace{0.2cm}

% ============================================================
% PPT 1：前向动态与计算图
% ============================================================
\subsection*{PPT 1：前向时空动态与不可微源头}
\paragraph{\textbf{PPT内容要点：}}
\begin{itemize}
    \item 时空参数：层索引 \(l\)，时间步 \(t\)。
    \item 神经元状态：膜电位 \(u_i^l[t]\)，脉冲 \(s_i^l[t]\)。
    \item LIF离散动力学（泄漏 + 硬重置）与 脉冲生成（阶跃函数，不可导）：
\begin{equation}
    \left\{
    \begin{aligned}
        &u_i\left[t + 1\right] = \lambda (u_i[t] - V_{th}s_i[t]) + \sum_j w_{ij}s_j[t] + b_i,\\
        &s_i[t + 1] = H(u_i\left[t+1\right] - V_{th}),\\
    \end{aligned}
    \right.
    \label{eq.discrete}
\end{equation}

\end{itemize}
\paragraph{\textbf{讲稿：}}
\begin{itemize}
    \item 训练时网络沿时间和空间双向展开，构成时空计算图。
    \item 每个节点由 \((u, s)\) 描述，更新遵循LIF方程。
    \item 方程含两项传播机制：泄漏项 \(\lambda u\)（记忆衰减）与重置项 \(-V_{th} s\)（硬复位）。
    \item 脉冲函数 \(H\) 几乎处处导数为0，此为BPTT需代理梯度的根本原因。
\end{itemize}

% ============================================================
% PPT 2：BPTT完整公式（核心）
% ============================================================
\subsection*{PPT 2：BPTT完整梯度公式（标准形式）}
\paragraph{\textbf{PPT内容要点：}}
\begin{itemize}
    \item 损失定义：
  \begin{equation}
    L[t]=\frac{1}{T}\mathcal{L}\left(\mathbf{s}^N[t], \mathbf{y}\right), L = \sum_{t=1}^T L[t]
    \label{instantaneous loss}
  \vspace{-2mm}
  \end{equation}
    \item 对权重 \(\mathbf{W}^l\) 的梯度（论文式(3)）：
    \[
    \boxed{
    \begin{aligned}
    \frac{\partial L}{\partial \mathbf{W}^l} 
    &= \sum_{t=1}^T 
    \frac{\partial L}{\partial \mathbf{s}^{l+1}[t]}
    \cdot
    {\color{red}{\frac{\partial \mathbf{s}^{l+1}[t]}{\partial \mathbf{u}^{l+1}[t]}}}
    \cdot
    \Bigg(
    \frac{\partial \mathbf{u}^{l+1}[t]}{\partial \mathbf{W}^l}
    \\
    &\quad + 
    \sum_{\tau < t} 
    \Bigg[
    \prod_{i=t-1}^{\tau}
    \left(
    {\color[rgb]{0,0.69,0.31}{\frac{\partial \mathbf{u}^{l+1}[i+1]}{\partial \mathbf{u}^{l+1}[i]}}}
    +
    {\color{cyan}{\frac{\partial \mathbf{u}^{l+1}[i+1]}{\partial \mathbf{s}^{l+1}[i]}}}
    \cdot
    {\color{red}{\frac{\partial \mathbf{s}^{l+1}[i]}{\partial \mathbf{u}^{l+1}[i]}}}
    \right)
    \Bigg]
    \cdot
    \frac{\partial \mathbf{u}^{l+1}[\tau]}{\partial \mathbf{W}^l}
    \Bigg)
    \end{aligned}
    }
    \]
\end{itemize}
\paragraph{\textbf{讲稿：}}
\begin{itemize}
    \item 公式总体大结构：外层对 \(t\) 求和，根据上述的\(L = \sum_{t=1}^T L[t]\), 内部由“当前的膜电位的贡献”与“历史回溯的膜电位的贡献”相加。
    \item 红色项 \({\color{red}{\frac{\partial s}{\partial u}}}\)：脉冲对膜电位的敏感度（不可导，后续用代理梯度替代）。
    \item 绿色项 \({\color[rgb]{0,0.69,0.31}{\frac{\partial u[i+1]}{\partial u[i]}}}\)：膜电位随着时间自动地进行衰减，等于 \(\lambda \mathbf{I}\)（泄漏路径）。
    \item 青色项 \({\color{cyan}{\frac{\partial u[i+1]}{\partial s[i]}}}\)：脉冲对下一时刻膜电位的影响，等于 \(-\lambda V_{th}\mathbf{I}\)（重置路径）。
    \item \(\color{cyan}{CyanPart}\times\color{red}{RedPart}+\color{green}{GreenPart}\)两项在连乘内相加，代表误差沿时间的两条并行传播通道。
\end{itemize}




% ============================================================
% PPT 3：梯度总路径与当前直连项（空间传播）
% ============================================================
\subsection*{PPT 3：$\frac{\partial L}{\partial \mathbf{W}^l}$ 的完整传播路径与“当前直连项”}

\paragraph{\textbf{PPT内容要点：}}
\begin{itemize}
    \item \textbf{计算目标}：损失 \(L\) 对第 \(l\) 层权重 \(\mathbf{W}^l\) 的梯度。该权重连接 \textbf{上游层 \(l\)（脉冲 \(\mathbf{s}^l\)）} 与 \textbf{下游层 \(l+1\)（膜电位 \(\mathbf{u}^{l+1}\)）}；
    \item \textbf{时间下标约定}：严格按公式，\(u^{l+1}[t+1]\) 由 \(s^l[t]\) 驱动。用 \(t\) 表示当前响应时刻。因此梯度公式中写为 \(\frac{\partial \mathbf{u}^{l+1}[t]}{\partial \mathbf{W}^l} = \mathbf{s}^{l}[t]\)；
    \item \textbf{外层求和} \(\sum_{t=1}^T\)：权重在所有时刻共享，需累加各时刻的梯度贡献。
    \item \textbf{误差传入项} \(\frac{\partial L}{\partial \mathbf{s}^{l+1}[t]}\)：由输出层通过空间反向传播（层间BP）递推到节点 \((l+1, t)\) 的脉冲端。
    \item \textbf{红色代理梯度} \({\color{red}{\frac{\partial \mathbf{s}}{\partial \mathbf{u}}}}\)：在同一节点 \((l+1, t)\) 内部，将误差从“脉冲”映射到“膜电位”。
    \item \textbf{当前直连项} \(\frac{\partial \mathbf{u}^{l+1}[t]}{\partial \mathbf{W}^l} = \mathbf{s}^{l}[t]\)：空间上的直接传播。由上游节点 \((l, t)\) 的脉冲，通过权重 \(\mathbf{W}^l\) 直接传到下游节点 \((l+1, t)\) 的膜电位。
    \item \textbf{无时间回溯}：当前项只依赖当前时刻输入 \(s^l[t]\)，不含 \(\prod\) 或 \(\sum_{\tau < t}\)，纯空间部分
\end{itemize}

\paragraph{\textbf{讲稿：}}
\begin{itemize}
    \item 整个梯度刻画的是：上游层 \(l\) 在所有时刻的脉冲，如何通过权重影响下游层 \(l+1\) 的膜电位，并最终改变损失。
    \item 为统一符号，论文将层间延迟吸收进时间索引，故 \(\frac{\partial u^{l+1}[t]}{\partial W^l} = s^l[t]\) 在约定下精确成立。
    \item 当前项的计算路径：节点 \((l, t)\) 脉冲 \(\xrightarrow{W^l}\) 节点 \((l+1, t)\) 膜电位。
    \item 反向传播时：高层误差 \(\frac{\partial L}{\partial s^{l+1}[t]}\) 到达节点 \((l+1, t)\) 脉冲端。
    \item 乘以红色代理梯度，将误差从脉冲传至同一节点的膜电位。
    \item 随后该膜电位误差乘以 \(\frac{\partial u^{l+1}[t]}{\partial s^l[t]} = W^l\)，沿空间反向传回上层节点 \((l, t)\) 的脉冲端。
    \item 当前项不含连乘或历史求和，因为膜电位对权重的直接导数只依赖当前输入脉冲，不依赖过去状态。
    \item 整体相乘因子为：误差项 × 代理梯度 × 当前输入脉冲。
\end{itemize}

% ============================================================
% PPT 4：历史回溯项（时间传播核）—— 详解相乘与相加
% ============================================================
\subsection*{PPT 4：公式拆解（核心）—— 时间传播核与“相乘、相加”的数学本质}

\paragraph{\textbf{PPT内容要点：}}
\begin{itemize}
    \item \textbf{历史回溯项结构}：
    \[
    \sum_{\tau < t} \left[ \prod_{i=t-1}^{\tau} \left( 
    {\color[rgb]{0,0.69,0.31}{\lambda \mathbf{I}}} 
    + {\color{cyan}{(-\lambda V_{th} \mathbf{I})}} \cdot {\color{red}{\sigma'}} 
    \right) \right] \cdot \frac{\partial u[\tau]}{\partial W}
    \]
    \item \textbf{节点传播路径（时间维度）}：同一层 \(l+1\) 内部，历史时刻 \(\tau\) 的膜电位节点 \(u[\tau]\)，通过神经元内部动态，一步一步传播到当前时刻 \(t\) 的膜电位节点 \(u[t]\)。
    \item \textbf{足迹项} \(\frac{\partial u[\tau]}{\partial W} = s[\tau]\)：历史时刻上游输入的脉冲。
    \item \textbf{相乘（\(\prod\)）：链式法则}：从 \(\tau\) 到 \(t\) 必须经过中间所有时刻 \(t-1, t-2, \dots, \tau\)，复合函数嵌套求导，根据链式法则，连续相乘。
    \item \textbf{括号内相加（\(+\)）：全微分}：同一个时间步 \(i \to i+1\) 膜电位更新包含两项：
    \[
    u[i+1] = \lambda u[i] - \lambda V_{th} s[i] + \cdots
    \]
    \(u[i+1]\) 对 \(u[i]\) 的依赖有两条并行路径：直接泄漏（绿色）和通过脉冲的间接重置（青色×红色）。两路径同时存在，总导数为二者之和（全微分公式：\(\frac{dz}{dx} = \frac{\partial z}{\partial x} + \frac{\partial z}{\partial y}\frac{dy}{dx}\)）。
    \item \textbf{外层求和（\(\sum_{\tau < t}\)）}：所有历史时刻 \(\tau = 1, \dots, t-1\) 都是独立的源头，各自通过时序核影响当前，总贡献为各源头之和。
    \item \textbf{青色项精确值（依据本论文公式(2)）}：
    \[
    {\color{cyan}{\frac{\partial u[i+1]}{\partial s[i]}}} = -\lambda V_{th} \mathbf{I}
    \]
    （对线性重置项 \(-\lambda V_{th} s[i]\) 直接求导）
\end{itemize}

\paragraph{\textbf{讲稿：}}
\begin{itemize}
    \item 历史回溯项计算的是：过去时刻 \(\tau\) 的输入脉冲，经过下游神经元内部时间演化，对当前时刻 \(t\) 膜电位的影响。
    \item 节点路径：同一层内部 \(u[\tau] \to u[\tau+1] \to \cdots \to u[t]\)。
    \item 为什么有连乘（\(\prod\)）？因为时间推进是顺序发生的，复合函数链式法则要求每一步的局部导数连续相乘。这就是“串联”关系。
    \item 为什么括号内是相加（\(+\)）？因为在相邻时间步 \(i \to i+1\) 中，膜电位受两种并行物理机制控制：泄漏（绿色）和重置（青色×红色）。二者互不排斥，同时作用。根据全微分原理，总传播系数等于各路径贡献的线性叠加。这就是“并联”关系。
    \item 青色项严格等于 \(-\lambda V_{th}\)。注意：这不是对阶跃函数求导，而是对线性项 \(-\lambda V_{th} s[i]\) 求导，自变量是 \(s\)，故直接得常数。
    \item 外层求和（\(\sum_{\tau < t}\)）是因为每个历史时刻都是独立的误差源头，总影响是各源头贡献的总和。
\end{itemize}

% ============================================================
% PPT 5：代理梯度、STBP定义及BPTT瓶颈
% ============================================================
\subsection*{PPT 5：代理梯度（SG）、STBP严格定义与BPTT固有瓶颈}

\paragraph{\textbf{PPT内容要点：}}
\begin{itemize}
    \item \textbf{代理梯度（Surrogate Gradient）的数学形式}：因 \(\frac{\partial s}{\partial u} = 0\)（a.e.），反向传播时用光滑函数导数替代。常用形式：
    \[
    \sigma'(x) = \frac{1}{a}\text{sign}\left(|x - V_{th}| < \frac{a}{2}\right) \quad \text{或} \quad \sigma'(x) = \frac{1}{a}\frac{e^{(V_{th}-x)/a}}{(1+e^{(V_{th}-x)/a})^2}.
    \]
    \item \textbf{使用范围}：代理梯度仅用于反向传播（红色项），前向传播仍为原始阶跃函数，不改变脉冲离散性。
    \item \textbf{STBP的严格学术定义}：
    \[
    \text{STBP (Spatio-Temporal Backpropagation)} \equiv \text{BPTT} + \text{Surrogate Gradient (SG)}.
    \]
    STBP并非独立新算法，而是BPTT在SNN场景下的工程化命名，其核心特征为完整的时间展开配合代理梯度。
    \item \textbf{BPTT的固有瓶颈}：
    \begin{itemize}
        \item \textbf{内存复杂度 \(\mathcal{O}(T)\)}：反向传播时必须存储全部时刻的膜电位 \(u[t]\) 和脉冲 \(s[t]\)，以计算时序核连乘。
        \item \textbf{非流式（Non-streaming）}：必须等整个长度为 \(T\) 的序列处理完毕后才能计算梯度并更新权重，无法在线学习。
        \item \textbf{无法适配神经形态硬件}：生物神经系统和Loihi等芯片要求逐时刻实时更新，BPTT的批处理特性与之矛盾。
    \end{itemize}
\end{itemize}

\paragraph{\textbf{讲稿：}}
\begin{itemize}
    \item 代理梯度是解决脉冲函数不可导的必要工程近似，仅反向使用，不影响前向计算。
    \item STBP在严格意义上就是BPTT配合代理梯度，二者数学等价，只是命名侧重点不同。
    \item BPTT的核心问题来源于时序核的连乘结构，必须存储全轨迹，内存随 \(T\) 线性增长。
    \item BPTT需整段序列处理完才能更新，属于离线批处理，与生物在线学习机制不兼容。
    \item OTTT通过丢弃重置耦合路径，将双向时序回溯简化为前向指数衰减累积，从而打破 \(\mathcal{O}(T)\) 内存限制并支持逐时刻更新。
    \item 理解标准BPTT的这些瓶颈，是理解OTTT贡献的绝对前提。
\end{itemize}


% ------------------------------------------------------------
% 总结
% ------------------------------------------------------------
\subsection*{总结}
\begin{itemize}
    \item BPTT完整保留泄漏与重置两条时间传播路径。
    \item 代理梯度唯一负责处理脉冲不可导性。
    \item 理解BPTT的每一项是理解OTTT改进的前提。
\end{itemize}





\section{OTTT 与 NDOT：针对 SNN 的在线训练方法}


% ============================================================
% 副标题列表（本部分内容概览）
% ============================================================
\subsection*{本部分核心结构}
\begin{itemize}[leftmargin=*, itemsep=0.3ex]
    \item[\textbf{1}] \textbf{BPTT 标准形式与核心瓶颈}：时间展开、代理梯度依赖、\(\mathcal{O}(T)\) 内存。
    \item[\textbf{2}] \textbf{OTTT 的近似解耦}：丢弃重置路径，时序核退化为固定 \(\lambda^{t-\tau}\)，实现 \(\mathcal{O}(1)\) 在线学习。
    \item[\textbf{3}] \textbf{NDOT 的物理重参数化}：从连续 LIF 微分方程导出动态比值 \(\mathbf{e}[t]\)，替代固定 \(\lambda\)。
    \item[\textbf{4}] \textbf{NDOT 完整推导（Appendix A.1）}：动态核 \(\mathbf{P}_{k,t}\) 的化简与递推公式的严格证明。
    \item[\textbf{5}] \textbf{本质区别汇总}：固定衰减 vs. 动态比值、近似 vs. 精确、变量定义对照表。
\end{itemize}
\vspace{0.3cm}

% ============================================================
% PPT 1：BPTT 标准形式与核心瓶颈
% ============================================================
\subsection*{PPT 1：BPTT 标准形式与核心瓶颈}

\paragraph{\textbf{PPT内容要点（含完整变量定义）：}}
\begin{itemize}
    \item \textbf{离散 LIF 动力学}：
    \[
    \mathbf{u}^{l}[t+1] = \lambda (\mathbf{u}^{l}[t] - V_{th}\mathbf{s}^{l}[t]) + \mathbf{W}^{l}\mathbf{s}^{l-1}[t+1]
    \]
    其中 \(\mathbf{u}^{l}[t]\) 为第 \(l\) 层 \(t\) 时刻膜电位，\(\mathbf{s}^{l}[t]\) 为脉冲，\(\lambda\) 为泄漏常数，\(V_{th}\) 为阈值，\(\mathbf{W}^{l}\) 为待训练权重。
    
    \item \textbf{BPTT 标准梯度公式}：
    \[
    \boxed{
    \begin{aligned}
    \frac{\partial L}{\partial \mathbf{W}^l}
    &=
    \sum_{t=1}^T
    \frac{\partial L}{\partial \mathbf{s}^{l+1}[t]}
    {\color{red}{\frac{\partial \mathbf{s}^{l+1}[t]}{\partial \mathbf{u}^{l+1}[t]}}}
    \Bigg(
    \frac{\partial \mathbf{u}^{l+1}[t]}{\partial \mathbf{W}^l}
    \\
    &\quad +
    \sum_{\tau < t}
    \prod_{i=t-1}^{\tau}
    \left(
    {\color[rgb]{0,0.69,0.31}{\frac{\partial \mathbf{u}^{l+1}[i+1]}{\partial \mathbf{u}^{l+1}[i]}}}
    +
    {\color{cyan}{\frac{\partial \mathbf{u}^{l+1}[i+1]}{\partial \mathbf{s}^{l+1}[i]}}}
    {\color{red}{\frac{\partial \mathbf{s}^{l+1}[i]}{\partial \mathbf{u}^{l+1}[i]}}}
    \right)
    \frac{\partial \mathbf{u}^{l+1}[\tau]}{\partial \mathbf{W}^l}
    \Bigg)
    \end{aligned}
    }
    \]
    
    \item \textbf{各项精确定义}：
    \begin{itemize}
        \item 红色项 \({\color{red}{\frac{\partial \mathbf{s}}{\partial \mathbf{u}}}}\)：脉冲发放函数对膜电位的导数（Heaviside 不可微，需代理梯度 SG）。
        \item 绿色项 \({\color[rgb]{0,0.69,0.31}{\frac{\partial \mathbf{u}[i+1]}{\partial \mathbf{u}[i]}}}\)：膜电位自传播，严格等于 \(\lambda \mathbf{I}\)（泄漏路径）。
        \item 青色项 \({\color{cyan}{\frac{\partial \mathbf{u}[i+1]}{\partial \mathbf{s}[i]}}}\)：脉冲对下一时刻膜电位的直接影响，严格等于 \(-\lambda V_{th} \mathbf{I}\)（重置路径）。
    \end{itemize}
    
    \item \textbf{BPTT 两大瓶颈}：
    \begin{enumerate}
        \item 时序连乘中因含 \({\color{red}{\frac{\partial s}{\partial u}}}\)，必须使用代理梯度（近似，理论不清晰）。
        \item 需存储全部 \(T\) 个时刻的 \(\mathbf{u}\) 和 \(\mathbf{s}\) 以计算连乘，内存复杂度 \(\mathcal{O}(T)\)，且无法在线更新。
    \end{enumerate}
\end{itemize}

\paragraph{\textbf{讲稿：}}
\begin{itemize}
    \item 离散 LIF 动力学由泄漏项、重置项、输入项三部分构成。
    \item BPTT 展开后，梯度包含“当前直接贡献”与“历史回溯贡献”两部分。
    \item 历史回溯中的连乘每步含两条路径：绿色（泄漏）和青色×红色（重置耦合）。
    \item 红色项因不可微必须用 SG，是 BPTT 近似性的根源。
    \item 连乘结构迫使反向传播时必须访问所有历史状态，导致内存随 \(T\) 线性增长。
\end{itemize}

% ============================================================
% PPT 2：OTTT 的近似解耦与递推实现
% ============================================================
\subsection*{PPT 2：OTTT 的近似解耦与递推实现}

\paragraph{\textbf{PPT内容要点（含完整推导）：}}
\begin{itemize}
    \item \textbf{核心近似}：在时序连乘中，不对红色项施加代理梯度，而是直接使用其真实导数（几乎处处为 0）：
    \[
    {\color{cyan}{\frac{\partial \mathbf{u}[i+1]}{\partial \mathbf{s}[i]}}}
    {\color{red}{\frac{\partial \mathbf{s}[i]}{\partial \mathbf{u}[i]}}}
    \approx 0
    \quad\Longrightarrow\quad
    \prod_{i=\tau}^{t-1}
    \left(
    {\color[rgb]{0,0.69,0.31}{\lambda \mathbf{I}}} + 0
    \right)
    = \lambda^{t-\tau} \mathbf{I}
    \]
    
    \item \textbf{简化后的梯度公式}\cite{xiao_online_2022}：
    \[
    \frac{\partial L}{\partial \mathbf{W}^l}
    =
    \sum_{t=1}^T
    \frac{\partial L}{\partial \mathbf{s}^{l+1}[t]}
    {\color{red}{\frac{\partial \mathbf{s}^{l+1}[t]}{\partial \mathbf{u}^{l+1}[t]}}}
    \left(
    \sum_{\tau \leq t} \lambda^{t-\tau}
    \frac{\partial \mathbf{u}^{l+1}[\tau]}{\partial \mathbf{W}^l}
    \right)
    \]
    
    \item \textbf{定义追踪变量 \(\hat{\mathbf{a}}^l[t]\)（突触前活动踪迹）}：
    \[
    \boxed{
    \hat{\mathbf{a}}^l[t] \triangleq \sum_{\tau=1}^{t} \lambda^{t-\tau} \mathbf{s}^l[\tau]
    }
    \]
    由于 \(\frac{\partial \mathbf{u}^{l+1}[\tau]}{\partial \mathbf{W}^l} = \mathbf{s}^l[\tau]\)（对权重的直接偏导等于输入脉冲），上式可重写为：
    \[
    \nabla_{\mathbf{W}^l}L = \sum_{t=1}^T \mathbf{g}_{\mathbf{u}^{l+1}}[t] \left(\hat{\mathbf{a}}^l[t]\right)^\top
    \]
    其中 \(\mathbf{g}_{\mathbf{u}^{l+1}}[t] = \left( \frac{\partial L}{\partial \mathbf{s}^{l+1}[t]} {\color{red}{\frac{\partial \mathbf{s}^{l+1}[t]}{\partial \mathbf{u}^{l+1}[t]}}} \right)^\top\) 为空间梯度。
    
    \item \textbf{前向递推（在线计算）}：
    \[
    \boxed{
    \hat{\mathbf{a}}^l[t] = \lambda \hat{\mathbf{a}}^l[t-1] + \mathbf{s}^l[t]
    }
    \]
    此式仅依赖上一时刻的 \(\hat{\mathbf{a}}\)，无需存储历史轨迹。
    
    \item \textbf{OTTT 简化效果}：
    \begin{itemize}
        \item 时序依赖退化为固定指数衰减 \(\lambda^{t-\tau}\)。
        \item 内存复杂度降至 \(\mathcal{O}(1)\)（与 \(T\) 无关）。
        \item 红色项仅在空间传播入口处保留（仍用 SG），时序回溯中完全消除。
    \end{itemize}
\end{itemize}

\paragraph{\textbf{讲稿：}}
\begin{itemize}
    \item OTTT 将连乘中的重置路径直接置零，因为其依赖的代理梯度不可靠。
    \item 时序核简化为纯泄漏 \(\lambda^{t-\tau}\)，成为固定常数衰减。
    \item 定义追踪变量 \(\hat{a}[t]\) 为历史输入脉冲的加权和，权重为 \(\lambda^{t-\tau}\)。
    \item 利用 \(\frac{\partial u[\tau]}{\partial W} = s[\tau]\)，梯度可分解为空间梯度 \(\mathbf{g}\) 与追踪变量 \(\hat{a}\) 的外积。
    \item \(\hat{a}\) 可前向递推计算，只需保留上一时刻值，实现常数内存。
    \item 但固定 \(\lambda\) 无法反映神经元状态的实际波动，这是 OTTT 的精度瓶颈。
\end{itemize}

% ============================================================
% PPT 3：NDOT 的物理重参数化 —— e[t] 的严格推导
% ============================================================
\subsection*{PPT 3：NDOT 的物理重参数化 —— \(\mathbf{e}[t]\) 的严格推导}

\paragraph{\textbf{PPT：}}
\begin{itemize}
    \item \textbf{OTTT 的局限}：固定 \(\lambda\) 是常数，但神经元膜电位的实际演化受充电、泄漏、重置共同影响，衰减率应动态变化。
    
    \item \textbf{NDOT 核心观点}：用连续 LIF 微分方程推导出的 \textit{e(t), representing the continuous temporal dependency}替代固定 \(\lambda\)。
    
    \item \textbf{Step 1：连续 LIF 方程}\cite{jiang_ndot_nodate}：
    \[
    u'(t) = -u(t) + I(t)
    \]
    其中 \(I(t)\) 为输入电流。
    
    \item \textbf{Step 2：定义连续时序依赖 \(e(t)\)}（隐函数关系）：
    \[
    e(t) \triangleq \frac{\partial u(t+1)}{\partial u(t)}
    = \frac{u'(t+1)}{u'(t)}
    \]
    \textit{推导依据}：由隐函数定理，\(u(t+1)\) 对 \(u(t)\) 的偏导等于两者对时间导数之比。
    
    \item \textbf{Step 3：代入 LIF 方程消去 \(I(t)\)}：
    \[
    e(t) = \frac{-u(t+1) + I(t+1)}{-u(t) + I(t)}
    = \frac{u(t+1) - I(t+1)}{u(t) - I(t)}
    \]
    
    \item \textbf{Step 4：离散化并代入离散 LIF 方程}：
    在离散时间步 \(t\)，\(I[t]\) 满足 \(u[t+1] = \lambda(u[t] - V_{th}s[t]) + I[t+1]\)。
    因此 \(u[t+1] - I[t+1] = \lambda(u[t] - V_{th}s[t])\)。
    同时由前一时刻方程 \(I[t] = u[t] - \lambda^{-1}u[t-1]\)，代入分母化简得：
    \[
    \boxed{
    \mathbf{e}^{l}[t] =
    \frac{\mathbf{u}^{l}[t] - V_{th}\mathbf{s}^{l}[t]}
         {\mathbf{u}^{l}[t-1] - V_{th}\mathbf{s}^{l}[t-1]}
    }
    \]
    \textit{注意}：此式完全不包含代理梯度，纯为前向可计算的动态比值。
    
    \item \textbf{变量定义}：\(\mathbf{e}^{l}[t]\) 表示第 \(l\) 层从时刻 \(t-1\) 到 \(t\) 的膜电位状态缩放因子，是 NDOT 的核心时序核。
\end{itemize}

\paragraph{\textbf{讲稿：}}
\begin{itemize}
    \item OTTT 的固定 \(\lambda\) 是常数，无法适应膜电位的实际波动。
    \item NDOT 从连续 LIF 微分方程出发，利用隐函数关系定义动态敏感度 \(e(t)\)。
    \item 通过代入连续方程消去输入电流，再离散化并代入离散 LIF 方程，最终得到闭式解 \(\mathbf{e}[t]\)。
    \item \(\mathbf{e}[t]\) 是相邻时刻“重置后膜电位”的比值，完全由前向状态决定。
    \item 该推导全程无需任何代理梯度，因此比 OTTT 的固定衰减更忠实于神经元物理。
\end{itemize}

% ============================================================
% PPT 4：NDOT 完整推导 —— P_{k,t}、化简与递推公式
% ============================================================
\subsection*{PPT 4：NDOT 完整梯度推导（严格依据 Appendix A.1）}

\paragraph{\textbf{PPT内容要点（含完整推导链）：}}
\begin{itemize}
    \item \textbf{Step 1：定义动态时序传播核 \(\mathbf{P}_{k,t}^{l}\)}（NDOT 论文式(9)）：
    \[
    \mathbf{P}_{k,t}^{l} \triangleq \prod_{i=k}^{t-1} \mathbf{e}^{l}[i]
    = \mathbf{e}^{l}[k] \odot \mathbf{e}^{l}[k+1] \odot \cdots \odot \mathbf{e}^{l}[t-1]
    \]
    其中 \(\odot\) 为逐元素乘法（保持维度匹配）。
    
    \item \textbf{Step 2：代入 \(\mathbf{e}[i]\) 并执行化简}：
    \[
    \begin{aligned}
    \mathbf{P}_{k,t}^{l}
    &=
    \frac{u[k] - V_{th}s[k]}{u[k-1] - V_{th}s[k-1]}
    \odot
    \frac{u[k+1] - V_{th}s[k+1]}{u[k] - V_{th}s[k]}
    \odot
    \cdots
    \odot
    \frac{u[t-1] - V_{th}s[t-1]}{u[t-2] - V_{th}s[t-2]}
    \\
    &=
    \frac{\mathbf{u}^{l}[t-1] - V_{th}\mathbf{s}^{l}[t-1]}
         {\mathbf{u}^{l}[k-1] - V_{th}\mathbf{s}^{l}[k-1]}
    \end{aligned}
    \]
    中间项全部抵消，仅剩首尾比值。
    
    \item \textbf{Step 3：定义时序追踪变量 \(\hat{\mathbf{a}}^{l-1}[t]\)}（NDOT 论文式(11)）：
    \[
    \boxed{
    \hat{\mathbf{a}}^{l-1}[t] \triangleq
    \mathbf{s}^{l-1}[t] + \sum_{k < t} \mathbf{P}_{k,t}^{l-1} \odot \mathbf{s}^{l-1}[k]
    }
    \]
    该变量累加所有历史时刻 \(k\) 的输入脉冲，经动态核 \(\mathbf{P}_{k,t}\) 缩放后求和。
    
    \item \textbf{Step 4：推导相邻时刻递推关系（Appendix A.1 核心证明）}：
    由定义展开 \(\hat{\mathbf{a}}[t-1]\)：
    \[
    \hat{\mathbf{a}}[t-1] = \mathbf{s}[t-1] + \sum_{k < t-1} \mathbf{P}_{k,t-1} \odot \mathbf{s}[k]
    \]
    利用 \(\mathbf{P}_{k,t} = \mathbf{e}[t-1] \odot \mathbf{P}_{k,t-1}\)（核递推性质）：
    \[
    \begin{aligned}
    \mathbf{e}[t-1] \odot \hat{\mathbf{a}}[t-1]
    &= \mathbf{e}[t-1] \odot \mathbf{s}[t-1]
       + \sum_{k < t-1} \mathbf{P}_{k,t} \odot \mathbf{s}[k] \\
    &= \mathbf{P}_{t-1,t} \odot \mathbf{s}[t-1] + \sum_{k < t-1} \mathbf{P}_{k,t} \odot \mathbf{s}[k] \\
    &= \sum_{k < t} \mathbf{P}_{k,t} \odot \mathbf{s}[k]
    \end{aligned}
    \]
    因此：
    \[
    \boxed{
    \hat{\mathbf{a}}^{l}[t] = \mathbf{e}^{l}[t-1] \odot \hat{\mathbf{a}}^{l}[t-1] + \mathbf{s}^{l}[t]
    }
    \]
    
    \item \textbf{Step 5：完整梯度公式}（NDOT 论文式(10)）：
    \[
    \boxed{
    \nabla_{\mathbf{W}^l} \mathcal{L}
    =
    \sum_{t=1}^{T}
    \mathbf{g}_{\mathbf{u}^l}[t]
    \left(
    \hat{\mathbf{a}}^{l-1}[t]
    \right)^{\top}
    }
    \]
    其中 \(\mathbf{g}_{\mathbf{u}^l}[t] = \left( \frac{\partial \mathcal{L}}{\partial \mathbf{s}^l[t]} {\color{red}{\frac{\partial \mathbf{s}^l[t]}{\partial \mathbf{u}^l[t]}}} \right)^{\top}\) 为空间梯度（仍含 SG）。
\end{itemize}

\paragraph{\textbf{讲稿：}}
\begin{itemize}
    \item NDOT 定义动态核 \(\mathbf{P}_{k,t}\) 为 \(\mathbf{e}\) 的连乘，替代 OTTT 的 \(\lambda^{t-\tau}\)。
    \item 因 \(\mathbf{e}\) 是分式比值，连乘形成效应，化简为首尾比值，极大简化计算。
    \item 定义追踪变量 \(\hat{a}[t]\) 为所有历史脉冲经动态核缩放后的总和。
    \item 利用 \(\mathbf{P}_{k,t}\) 的递推性质，严格推导出 \(\hat{a}[t] = \mathbf{e}[t-1] \odot \hat{a}[t-1] + s[t]\)。
    \item 该递推式与 OTTT 形式一致，但核心系数从固定 \(\lambda\) 换为动态 \(\mathbf{e}[t-1]\)。
    \item 最终梯度分解为空间梯度 \(\mathbf{g}\) 与追踪变量 \(\hat{a}\) 的外积，实现 \(\mathcal{O}(1)\) 在线计算。
\end{itemize}

% ============================================================
% PPT 5：本质区别汇总
% ============================================================
\subsection*{PPT 5：OTTT 与 NDOT 本质区别汇总}

\paragraph{\textbf{PPT内容要点（单条目对照）：}}
\begin{itemize}
    \item \textbf{时序核（单步）}：
    \begin{itemize}
        \item OTTT：\(\epsilon[i] = \lambda\)，固定常数。
        \item NDOT：\(\mathbf{e}[i] = \dfrac{u[i] - V_{th}s[i]}{u[i-1] - V_{th}s[i-1]}\)，动态状态比值。
    \end{itemize}
    
    \item \textbf{历史传播核}：
    \begin{itemize}
        \item OTTT：\(\lambda^{t-k}\)，固定指数衰减。
        \item NDOT：\(\mathbf{P}_{k,t} = \prod_{i=k}^{t-1} \mathbf{e}[i] = \dfrac{u[t-1] - V_{th}s[t-1]}{u[k-1] - V_{th}s[k-1]}\)，动态化简。
    \end{itemize}
    
    \item \textbf{追踪变量递推}：
    \begin{itemize}
        \item OTTT：\(\hat{a}[t] = \lambda \hat{a}[t-1] + s[t]\)。
        \item NDOT：\(\hat{a}[t] = \mathbf{e}[t-1] \odot \hat{a}[t-1] + s[t]\)。
    \end{itemize}
    
    \item \textbf{额外存储}：
    \begin{itemize}
        \item OTTT：无。
        \item NDOT：需存 \(u[t] - V_{th}s[t]\) 以计算 \(\mathbf{e}[t]\)。
    \end{itemize}
    
    \item \textbf{方法论本质}：
    \begin{itemize}
        \item OTTT：对 BPTT 做\textbf{近似删除}（丢弃重置路径），时序退化为固定衰减。
        \item NDOT：对 BPTT 做\textbf{物理重参数化}（用连续动力学比值 \(e\) 替代离散近似 \(\epsilon\)），时序核变为动态状态比值。
    \end{itemize}
\end{itemize}

\paragraph{\textbf{共同基础：}}
\begin{itemize}
    \item 均实现 \(\mathcal{O}(1)\) 内存与在线更新。
    \item 梯度最终形式均为 \(\nabla L[t] = \mathbf{g}[t](\hat{a}[t])^\top\)。
    \item 空间传播入口处均保留红色项 \({\color{red}{\frac{\partial s}{\partial u}}}\)，使用代理梯度。
\end{itemize}

\paragraph{\textbf{讲稿（提纲式）：}}
\begin{itemize}
    \item OTTT 和 NDOT 都解决了 BPTT 的内存瓶颈，但路径不同。
    \item OTTT 删除重置路径，用固定 \(\lambda\) 做指数衰减，简单但损失时序精度。
    \item NDOT 用连续物理方程导出的动态 \(\mathbf{e}[t]\) 替代 \(\lambda\)，保留完整时序信息，且利用望远镜化简不增加计算负担。
    \item 递推形式完全一致，仅固定系数换为动态系数。
    \item NDOT 是 OTTT 的精度升级版，相同内存成本，更准确的梯度估计。
    \item 核心差异一句话：OTTT 是“删近似”，NDOT 是“用精确物理量替代近似”。
\end{itemize}




\section{S-TLLR 与 TESS：时间本地化与空间本地化学习}



% 参考文献：
% S-TLLR: \cite{apolinario_s-tllr_2024}
% TESS: \cite{apolinario_tess_2025}

% ============================================================
% 副标题列表
% ============================================================
\subsection*{本部分核心结构}
\begin{itemize}[leftmargin=*, itemsep=0.3ex]
    \item[\textbf{1}] \textbf{背景：BPTT 的瓶颈与三因子学习规则的机遇}
    \item[\textbf{2}] \textbf{S-TLLR：时间本地化学习规则}——广义 STDP 资格迹、次级激活函数 \(\Psi\)、因果/非因果融合
    \item[\textbf{3}] \textbf{TESS：时空全本地化学习规则}——从 \(\text{tr}\) 到 \(\mathbf{q}/\mathbf{h}\) 的实例化、本地学习信号生成（LSG）
    \item[\textbf{4}] \textbf{资格迹的进化：从抽象 \(\text{tr}\) 到显式张量 \(\mathbf{q}\) 与 \(\mathbf{h}\)}
    \item[\textbf{5}] \textbf{核心算法对比总结}——公式对照、复杂度差异、本地化维度进化
\end{itemize}
\vspace{0.3cm}

% ============================================================
% PPT 1：背景
% ============================================================
\subsection*{PPT 1：背景——BPTT 的瓶颈与三因子学习规则的机遇}

\paragraph{\textbf{PPT内容要点：}}
\begin{itemize}
    \item \textbf{BPTT 的计算瓶颈}：
    \[
    \frac{d\mathcal{L}}{d\mW^{(l)}} = \sum_{t=1}^T \frac{\partial \mathcal{L}}{\partial \vu^{(l)}[t]} \frac{\partial \vu^{(l)}[t]}{\partial \mW^{(l)}}
    \]
    内存复杂度 \(O(TLn)\)，时间复杂度 \(O(TLn^2)\)，随 \(T\) 线性增长
    
    \item \textbf{三因子学习规则（3F）的一般形式}：
    \[
    \Delta w_{ij} = \sum_t \delta_i[t] \cdot e_{ij}[t]
    \]
    其中 \(e_{ij}[t]\) 为资格迹 编码突触前/后活动历史，这个部分为时间信号
    
    \(\delta_i[t]\) 为学习信号 解决空间信用分配
    
    \item \textbf{标准资格迹的瓶颈}：
    \[
    e_{ij}[t] = \beta e_{ij}[t-1] + f(y_i[t])g(x_j[t])
    \]
    内存复杂度 \(O(n^2)\)，不适用于深度卷积 SNN。
    
    \item \textbf{STDP 的启示}：同时利用因果（前→后）和非因果（后→前）时序关系
\end{itemize}

\paragraph{\textbf{讲稿（提纲式）：}}
\begin{itemize}
    \item BPTT 的时间展开机制导致内存和计算成本随 \(T\) 线性增长，不适合边缘设备。
    \item 三因子学习规则提供生物学合理的替代方案，但标准资格迹的内存复杂度为 \(O(n^2)\)。
    \item STDP 同时利用因果和非因果时序关系，这启发了 S-TLLR 的设计。
    \item S-TLLR 和 TESS 分别从“时间本地”和“时空全本地”两个层面解决上述问题。
\end{itemize}

% ============================================================
% PPT 2：S-TLLR —— 时间本地化学习规则（含完整推导）
% ============================================================
\subsection*{PPT 2：S-TLLR 的完整推导与时间本地化实现}

\paragraph{\textbf{PPT内容要点（S-TLLR 核心算法公式）：}}
\begin{itemize}
    \item \textbf{广义 STDP 资格迹}\cite{apolinario_s-tllr_2024}：
    \[
    e_{ij}[t] = {\color{myred}\alpha_{pre}\Psi(u_i[t])\sum_{t'=0}^{t}\lambda_{pre}^{t-t'} x_j[t']}
    + {\color{myblue}\alpha_{post} x_j[t]\sum_{t'=0}^{t-1}\lambda_{post}^{t-t'} \Psi(u_i[t'])}
    \]
    \begin{itemize}
        \item {\color{myred}红色项（因果）}：当前后突触活动 \(\Psi(u_i[t])\) × 前突触历史踪迹
        \item {\color{myblue}蓝色项（非因果）}：当前前突触脉冲 \(x_j[t]\) × 后突触历史踪迹
    \end{itemize}
    
    \item \textbf{次级激活函数 \(\Psi\) 的作用}（论文附录）：
    \begin{itemize}
        \item \(\Psi\) 是一个以阈值为中心的钟形/三角形函数。
        \item 衡量“神经元离发放脉冲还有多远”，即使神经元未发放也能产生非零信号。
        \item 典型形式：三角形 \(\Psi(u) = 0.3 \times \max(1 - |u - v_{th}|, 0)\)。
        \item \textbf{与代理梯度（SG）的区别}：SG 用于反向传播（解决 \(\frac{\partial s}{\partial u}\) 不可导），\(\Psi\) 用于前向资格迹计算（衡量膜电位接近阈值的程度）。
    \end{itemize}
    
    \item \textbf{前向踪迹递推（时间本地化关键）}：
    \[
    \text{tr}(x_j)[t] = \lambda_{pre}\text{tr}(x_j)[t-1] + x_j[t]
    \]
    \[
    \text{tr}(\Psi(u_i))[t] = \lambda_{post}\text{tr}(\Psi(u_i))[t-1] + \Psi(u_i[t-1])
    \]
    
    \item \textbf{前向形式的资格迹}\cite{apolinario_s-tllr_2024}）：
    \[
    e_{ij}[t] = \alpha_{pre}\Psi(u_i[t])\text{tr}(x_j)[t] + \alpha_{post}x_j[t]\left(\text{tr}(\Psi(u_i))[t] - \Psi(u_i[t])\right)
    \]
    
    \item \textbf{学习信号与权重更新}\cite{apolinario_s-tllr_2024}：
    \[
    \delta_i^{(l)}[t] = 
    \begin{cases}
        \frac{\partial \mathcal{L}(\mathbf{y}^L[t], \mathbf{y}^*)}{\partial y_i^{(l)}[t]} & t \ge T_l \\
        0 & \text{otherwise}
    \end{cases}
    \]
    \[
    w_{ij} := w_{ij} + \rho \sum_{t=T_l}^{T} \delta_i[t] e_{ij}[t]
    \]
    （学习信号通过 BP 或 DFA 生成——时间本地，但空间非本地）
    
    \item \textbf{复杂度}：内存 \(O(Ln)\)，时间 \(O(Ln^2)\)。
\end{itemize}

\paragraph{\textbf{讲稿（提纲式）：}}
\begin{itemize}
    \item S-TLLR 的资格迹包含两项：因果项（红色，前→后）和非因果项（蓝色，后→前）。
    \item 次级激活函数 \(\Psi\) 与代理梯度不同——它在前向计算中衡量膜电位的“发放潜力”，即使神经元沉默也能产生学习信号。
    \item 通过踪迹变量的前向递推，资格迹可在当前时刻完全计算，无需存储历史状态。
    \item 学习信号 \(\delta\) 仍来自 BP 或 DFA，实现时间本地但空间非本地。
    \item \(\alpha_{post}\) 控制非因果项：\(-1\) 惩罚非因果，\(+1\) 奖励同步性，\(0\) 排除。
    \item 内存复杂度 \(O(Ln)\)，证明了“时间本地学习”的可行性。
\end{itemize}

% ============================================================
% PPT 3：TESS —— 从 tr 到 q/h 的实例化与空间本地化
% ============================================================
\subsection*{PPT 3：TESS 的时空全本地化实现与资格迹进化}

\paragraph{\textbf{PPT内容要点（TESS 核心算法公式）：}}
\begin{itemize}
    \item \textbf{使用显式张量 \(\mathbf{q}\) 与 \(\mathbf{h}\)}
    \begin{itemize}
        \item S-TLLR 的 \(\text{tr}(x_j)\) \(\rightarrow\)TESS 的 \(\mathbf{q}^{(l)}[t]\)（前突触踪迹张量）
        \item S-TLLR 的 \(\text{tr}(\Psi(u_i))\)） \(\rightarrow\) TESS 的 \(\mathbf{h}^{(l)}[t]\)（后突触踪迹张量）
    \end{itemize}
    
    \item \textbf{前突触踪迹 \(\mathbf{q}\)}（\cite{apolinario_tess_2025}(7)）：
    \[
    {\color{myred}\mathbf{q}^{(l)}[t] = \lambda_{pre}\mathbf{q}^{(l)}[t-1] + \mathbf{o}^{(l-1)}[t]}
    \]
    存储输入历史，维度 \(n^{(l-1)}\)。
    
    \item \textbf{后突触踪迹 \(\mathbf{h}\)}（\cite{apolinario_tess_2025}(9)）：
    \[
    {\color{myblue}\mathbf{h}^{(l)}[t] = \lambda_{post}\mathbf{h}^{(l)}[t-1] + \Psi(\mathbf{u}^{(l)}[t-1])}
    \]
    存储膜电位活跃历史，维度 \(n^{(l)}\)。
    
    \item \textbf{资格迹的向量化形式}（\cite{apolinario_tess_2025}(8)(10)）：
    \[
    {\color{myred}\mathbf{e}_{pre}^{(l)}[t] = \alpha_{pre}\Psi(\mathbf{u}^{(l)}[t]) \otimes \mathbf{q}^{(l)}[t]}
    \]
    \[
    {\color{myblue}\mathbf{e}_{post}^{(l)}[t] = \alpha_{post}\mathbf{h}^{(l)}[t] \otimes \mathbf{o}^{(l-1)}[t]}
    \]
    
    \item \textbf{进化 2：学习信号的本地生成——LSG 机制}（\cite{apolinario_tess_2025}(11)）：
    \[
    {\color{mygreen}\mathbf{v}_m^{(l)}[t] = \mathbf{B}^{(l)\top} \left( f(\mathbf{B}^{(l)} \mathbf{o}^{(l)}[t]) - \mathbf{y}^* \right)}
    \]
    \begin{itemize}
        \item \(\mathbf{B}^{(l)}\)：固定二进制矩阵（列向量为方波函数，准正交）
        \item \(f(\cdot)\)：softmax（分类）或 identity（回归）
        \item 完全在层内生成，不依赖其他层状态 → \textbf{空间本地化}
    \end{itemize}
    
    \item \textbf{完整权重更新}（\cite{apolinario_tess_2025}(12)-(14)）：
    \[
    \Delta \mathbf{W}_{pre}^{(l)}[t] = \left( {\color{mygreen}\mathbf{v}_m^{(l)}[t] \odot} {\color{myred}\alpha_{pre}\Psi(\mathbf{u}^{(l)}[t])} \right) \otimes \mathbf{q}^{(l)}[t]
    \]
    \[
    \Delta \mathbf{W}_{post}^{(l)}[t] = \left( {\color{mygreen}\mathbf{v}_m^{(l)}[t] \odot} {\color{myblue}\alpha_{post}\mathbf{h}^{(l)}[t]} \right) \otimes \mathbf{o}^{(l-1)}[t]
    \]
    \[
    \Delta \mathbf{W}^{(l)}[t] = \Delta \mathbf{W}_{pre}^{(l)}[t] + \Delta \mathbf{W}_{post}^{(l)}[t]
    \]
    
    \item \textbf{复杂度}：内存 \(O(Ln)\)，时间 \(O(LCn)\)（\(C\) 为类别数，\(C \ll n\)）。
\end{itemize}

\paragraph{\textbf{讲稿（提纲式）：}}
\begin{itemize}
    \item TESS 将 S-TLLR 中抽象的 \(\text{tr}\) 函数实例化为两个显式张量：\(\mathbf{q}\)（前突触踪迹）和 \(\mathbf{h}\)（后突触踪迹）。
    \item 这一实例化使算法能够以标准的张量运算高效运行于 VGG、ResNet 等深度架构。
    \item 更根本的进化是 LSG 机制：用固定矩阵 \(\mathbf{B}\) 直接在当前层输出上生成学习信号，彻底切断层间依赖。
    \item \(\mathbf{B}\) 的列向量为方波函数，准正交且硬件友好。
    \item TESS 实现了真正的“时空全本地”学习——时间和空间信用分配均只依赖本地信息。
    \item 时间复杂度从 S-TLLR 的 \(O(Ln^2)\) 降至 \(O(LCn)\)。
\end{itemize}

% ============================================================
% PPT 4：S-TLLR 与 TESS 的本质区别 —— 公式对照与复杂度对比
% ============================================================
\subsection*{PPT 4：S-TLLR 与 TESS 的本质区别汇总}

\paragraph{\textbf{PPT内容要点（核心算法对比）：}}

\textbf{区别 1：学习信号的生成方式（最关键）}
\begin{itemize}
    \item \textbf{S-TLLR}（论文式(10)）：
    \[
    \delta_i^{(l)}[t] = \frac{\partial \mathcal{L}}{\partial y_i^{(l)}[t]} \quad \text{（依赖 BP/DFA，空间非本地）}
    \]
    \item \textbf{TESS}（\cite{apolinario_tess_2025}(11)）：
    \[
    \mathbf{v}_m^{(l)}[t] = \mathbf{B}^{(l)\top} \left( f(\mathbf{B}^{(l)} \mathbf{o}^{(l)}[t]) - \mathbf{y}^* \right) \quad \text{（完全层内生成，空间本地）}
    \]
\end{itemize}

\textbf{区别 2：资格迹的形式（抽象 vs. 实例化）}
\begin{itemize}
    \item \textbf{S-TLLR}（标量，\cite{apolinario_s-tllr_2024}(9)）：
    \[
    e_{ij}[t] = \alpha_{pre}\Psi(u_i[t])\text{tr}(x_j)[t] + \alpha_{post}x_j[t](\text{tr}(\Psi(u_i))[t] - \Psi(u_i[t]))
    \]
    \item \textbf{TESS}（向量化，\cite{apolinario_tess_2025}(8)(10)）：
    \[
    \mathbf{e}_{pre}^{(l)}[t] = \alpha_{pre}\Psi(\mathbf{u}^{(l)}[t]) \otimes \mathbf{q}^{(l)}[t], \quad
    \mathbf{e}_{post}^{(l)}[t] = \alpha_{post}\mathbf{h}^{(l)}[t] \otimes \mathbf{o}^{(l-1)}[t]
    \]
\end{itemize}

\textbf{区别 3：时间复杂度}
\begin{itemize}
    \item S-TLLR：\(\text{MAC}_{\text{S-TLLR}} = 3(T-T_l) \sum n^{(l)} n^{(l-1)} = O(Ln^2)\)
    \item TESS：\(\text{MAC}_{\text{TESS}} = (T-t_l) \sum 2 \times n^{(l)} \times C = O(LCn)\)
\end{itemize}

\textbf{区别 4：本地化维度}
\begin{itemize}
    \item S-TLLR：{\color{myred}仅时间本地}（时间信用分配本地化，空间依赖 BP/DFA）
    \item TESS：{\color{mygreen}时空全本地}（时间 + 空间信用分配均本地化）
\end{itemize}

\paragraph{\textbf{讲稿（提纲式）：}}
\begin{itemize}
    \item S-TLLR 和 TESS 的核心区别在于学习信号的生成方式。
    \item S-TLLR 使用 BP/DFA 生成 \(\delta\)，依赖全局误差信号；TESS 使用 LSG 生成 \(\mathbf{v}_m\)，完全本地。
    \item 资格迹形式上，S-TLLR 为标量抽象，TESS 为向量化张量。
    \item 时间复杂度上，TESS 通过 LSG 将 \(O(Ln^2)\) 降至 \(O(LCn)\)。
    \item TESS 是 S-TLLR 的自然进化，从“时间本地”到“时空全本地”。
\end{itemize}

% ============================================================
% PPT 5：总结 —— 完整的算法进化路线图
% ============================================================
\subsection*{PPT 5：总结 —— 从 S-TLLR 到 TESS 的进化路线图}

\paragraph{\textbf{PPT内容要点：}}
\begin{itemize}
    \item \textbf{共同基础}：
    \begin{itemize}
        \item 均基于三因子学习规则（资格迹 + 学习信号）
        \item 均受 STDP 启发，融合因果与非因果时序关系
        \item 均实现 \(O(Ln)\) 内存复杂度，与时间步数 \(T\) 无关
        \item 均使用次级激活函数 \(\Psi\)（但角色略有不同）
    \end{itemize}
    
    \item \textbf{关键进化}：
    \begin{itemize}
        \item 学习信号：BP/DFA（S-TLLR）→ LSG（TESS）——从“误差分发”到“目标投射”
        \item 踪迹表示：抽象 \(\text{tr}\)（S-TLLR）→ 显式张量 \(\mathbf{q}/\mathbf{h}\)（TESS）
        \item 复杂度：\(O(Ln^2)\)（S-TLLR）→ \(O(LCn)\)（TESS）
        \item 本地化维度：仅时间本地 → 时空全本地
    \end{itemize}
    
    \item \textbf{方法论本质}：
    \begin{quote}
        S-TLLR 证明了“时间本地学习”是可行的；TESS 在此基础上用 LSG 替换 BP，实现了“时空全本地学习”，在不增加内存的前提下大幅降低计算复杂度。
    \end{quote}
    
    \item \textbf{一句话总结}：
    \begin{quote}
        S-TLLR 是将 STDP 生物机制融入时间维度的尝试；TESS 通过固定随机投影将误差信号也彻底本地化，实现了完全意义上的时空本地学习，在保持接近 BPTT 精度的同时将计算成本降低 205–661 倍。
    \end{quote}
\end{itemize}

\paragraph{\textbf{讲稿（提纲式）：}}
\begin{itemize}
    \item 两篇工作共同构建了从“时间本地”到“时空全本地”的完整技术路线。
    \item S-TLLR 证明时间本地学习的可行性；TESS 证明全本地学习的可行性。
    \item 核心进化：LSG 机制用固定随机投影替代 BP，切断了层间依赖。
    \item 两者均保持 \(O(Ln)\) 内存，但 TESS 将时间复杂度降至 \(O(LCn)\)。
    \item 这是一个从“误差分发”到“目标投射”的方法论飞跃。
    \item 两者共同为边缘设备上的高效 SNN 训练提供了理论基础与算法支撑。
\end{itemize}





\bibliographystyle{IEEEtran}
\bibliography{SNN_OnlineLearningSummary}
\end{document}
